\documentclass[10pt,a4paper]{amsart}
\usepackage[margin=19mm]{geometry}
\usepackage{amsmath,amssymb,mathtools}
\usepackage[scr=rsfs,cal=euler]{mathalfa}
\usepackage{xfrac}
\usepackage[unicode,hidelinks,bookmarks=false]{hyperref}
\newcommand{\ie}{i.e.\ }
\newcommand{\cf}{cf.\ }
\newcommand{\resp}{respectively\ }
\newcommand{\Subsection}{\subsection}
\providecommand{\nobreakdash}{\nobreak\hbox{-}}
\setlength{\emergencystretch}{2em}
\multlinegap0pt
\usepackage{bm}
\usepackage{amssymb}
\usepackage{tikz}
\usepackage{siunitx}
\usepackage[shortlabels]{enumitem}
\newtheorem{lemma}{Lemma}
\title{Infinite-dimensional spherical-radial decomposition for estimation of probabilistic constraints}
\author{Kewei Wang, Georg Stadler}
\date{Source: arXiv:2603.19907v1 (CC BY 4.0)}
\begin{document}
\maketitle

\noindent\emph{Source and licence.} Infinite-dimensional spherical-radial decomposition for estimation of probabilistic constraints. Version arXiv:2603.19907v1 (CC BY 4.0), distributed by arXiv under the Creative Commons Attribution 4.0 International licence, http://creativecommons.org/licenses/by/4.0/, as recorded in the arXiv licence metadata for this exact version and shown on the abstract page https://arxiv.org/abs/2603.19907v1. Full licence notice: \texttt{licenses/arxiv-2603.19907-CC-BY-4.0.txt}.\par\smallskip

\noindent\textit{Editorial scope.} Counted unit: the article's lemma Eq:phi\_u\_dis2 (source0.tex lines 595--600). The article's complete original proof of it is delivered with it (source0.tex lines 602--618, one \texttt{proof} environment of 17 lines). No mathematics is written, repaired or re-derived: the statement and the proof are the article's own lines, in the article's own notation, and no retained line is substituted or re-flowed.  Retained uncounted as the source-local context the proof needs: lines 513--515.  Nothing was omitted from any retained span, so the package carries no omission record.  No retained line contains a citation, so the wrapper carries no bibliography and no bibliography entry.  No retained line contains a figure environment, an image-inclusion command or drawing code, so no image is delivered and the package declares no asset dependency.  Editorial formatting only, in addition: an AMS \texttt{amsart} preamble that restates the article's own declarations and macros used by the retained lines, namely the environments \texttt{lemma} on separate counters, together with the standard packages those lines need.  The retained environments print in this document as Lemma 1 in order of appearance, whereas the article prints the same items with the numbering of its own full document.  The complete native source of the article as distributed in the arXiv e-print for this exact version is bundled unchanged as \texttt{sources/196812/source0.tex}.
\begin{equation}\label{Eq:phi_u_dis}
    \varphi(u)\approx \tilde{\varphi}_N(u):=\frac{1}{N}\sum_{i=1}^N \left[ \max{\{0, F_{\chi_K}(\rho_{\mathrm{out}}(u,\boldsymbol{v}_i,z_i))-F_{\chi_K}(\rho_{\mathrm{in}}(u,\boldsymbol{v}_i,z_i))\}} \right].
\end{equation}

\begin{lemma}
    In the setting and with the above definitions, it holds that:
    \begin{equation}\label{Eq:phi_u_dis2}
        \tilde{\varphi}_N(u)=\frac{1}{N}\sum_{i=1}^N \left[ \max{\{0, F_{\chi_K}(\tilde{\rho}_{\mathrm{out}}(u,\boldsymbol{v}_i,z_i))-F_{\chi_K}(\tilde{\rho}_{\mathrm{in}}(u,\boldsymbol{v}_i,z_i))\}} \right].
    \end{equation}
\end{lemma}

\begin{proof}
    For each constraint $j \in \{1, \dots, M\}$, we have
    \begin{equation*}
        I_j := \{r \ge 0 : g_j(u,\bar{\xi}+z+r L_K\boldsymbol{v}) \le 0\} = [\rho_{j,\mathrm{in}}, \rho_{j,\mathrm{out}}].
    \end{equation*}
    Specifically, for $j \in J_{\mathrm{i}}$, the center $\bar{\xi}+z$ is inside the feasible half-space. Therefore, $\rho_{j,\mathrm{in}} = 0$ and $I_j = [0, \rho_{j,\mathrm{out}}]$. For $j \in J_{\mathrm{o}}$, the center is outside, and the ray never leaves the feasible half-space after intersecting it. Therefore, $\rho_{j,\mathrm{out}} = \infty$ and $I_j = [\rho_{j,\mathrm{in}}, \infty)$.
    
    The overall feasible set along the ray is the intersection of these intervals, given by the (possibly empty) interval:
    \begin{equation*}
        \begin{aligned}
            \left\{r\geq 0: g(u,\bar{\xi}+z+r L_K\boldsymbol{v})\leq 0\right\} &= \bigcap_{j=1}^M I_j = %\Big[ \max_{j \in J_{\mathrm{o}}} \rho_{j,\mathrm{in}}, \min_{j \in J_{\mathrm{i}}} \rho_{j,\mathrm{out}} \Big] \\
            %&=
            [\tilde{\rho}_{\mathrm{in}}(u,\boldsymbol{v},z), \tilde{\rho}_{\mathrm{out}}(u,\boldsymbol{v},z)].
        \end{aligned}
    \end{equation*}
  If this intersection is empty, i.e., $\tilde{\rho}_{\mathrm{in}} > \tilde{\rho}_{\mathrm{out}}$, then the operator $\max\{0,\cdot\}$ in~\eqref{Eq:phi_u_dis2} correctly sets the negative CDF difference to $0$. Consequently, applying the same reasoning used to obtain~\eqref{Eq:phi_u_dis}, $\tilde{\varphi}_N(u)$ can be equivalently expressed as in~\eqref{Eq:phi_u_dis2}.
\end{proof}

\end{document}
